\chapter{Conclusion and Future Scope}
Development of low-cost and efficient adsorbents for fluoride removal is important for water treatment applications. In the present work, coconut fiber was selected as a biomass adsorbent and surface modification using Glycidyl Trimethyl Ammonium Chloride (GTMAC) was carried out to improve adsorption characteristics. DOE using Central Composite Design (CCD) was used for planning adsorption studies, and preliminary Langmuir adsorption model diagnostics were analyzed during Phase 2 of the work. This chapter presents the conclusions, future scope, and learning outcomes of the project. 

\section{Conclusion}
The present work focused on the preparation and GTMAC functionalization of coconut fiber for fluoride adsorption studies. Biomaterial preparation was carried out through drying, grinding, and sieving processes to obtain uniform particle size suitable for adsorption studies. 

CCD experimental runs were generated using Design-Expert software by varying pH, adsorbent dose, fluoride concentration, and contact time. Preliminary Langmuir model diagnostics showed acceptable prediction performance with an R² value of 0.8158 and RMSE value of 0.5278 mg/g. 

The prepared GTMAC-modified coconut fiber and CCD experimental design provide a good foundation for future fluoride adsorption studies and optimization analysis. 

\section{Future Scope}
Future work includes carrying out complete batch adsorption studies and fluoride removal experiments using the prepared GTMAC-modified coconut fiber. SEM and FTIR analysis can be performed to study surface morphology and confirm successful GTMAC functionalization.
\section{Learning Outcomes of the Project}
\begin{itemize}
\item Knowledge of GTMAC surface functionalization.  
\item Understanding of DOE and CCD experimental design.  
\item Knowledge of Langmuir adsorption modelling.  
\item Exposure to FTIR and SEM characterization techniques.   
\item Understanding of adsorption data analysis using Design-Expert software.  

\end{itemize}

